\section*{Chapitre \ref{ch:b2:anthropogenic-climate-change} --- Changement climatique anthropique}

\begin{solution}{exo:b2:anthropogenic-climate-change:1}
(a) $3.75 - 0.60 + 0.22 + 0.45 + 0.11 = \qty{3.93}{W.m^{-2}}$, les ajustements ajoutant
\qty{0.18}{W.m^{-2}}. (b) $3.93/3.71 = 1.059$~: \qty{5.9}{\percent} de plus. (c) L'ajustement de
température de la troposphère, $-0.60$~: une troposphère plus chaude émet davantage d'infrarouge vers
l'espace, ce qui retire une partie de l'énergie supplémentaire avant que la surface n'ait réagi.
\end{solution}

\begin{solution}{exo:b2:anthropogenic-climate-change:2}
(a) $4\pi(\num{6.371e6})^2 = \qty{5.10e14}{m^2}$. (b) Sur 47 ans, $\num{434.9e21}/(47 \times
\num{3.156e7} \times \num{5.10e14}) = \qty{0.575}{W.m^{-2}}$~; sur 12 ans, $\num{152.4e21}/(12
\times \num{3.156e7} \times \num{5.10e14}) = \qty{0.789}{W.m^{-2}}$, soit les 0.57 et 0.79 de
l'évaluation. (c) $0.91 \times 434.9 = \qty{396}{ZJ}$.
\end{solution}

\begin{solution}{exo:b2:anthropogenic-climate-change:3}
(a) $-3.22 + 1.30 + 0.35 + 0.42 - 0.01 = \num{-1.16}$, la valeur nette. (b) $3.93/1.16 =
\qty{3.39}{\celsius}$~; $3.93/3.22 = \qty{1.22}{\celsius}$. (c) $3.39/1.22 = 2.78$~: les
rétroactions multiplient presque par trois le réchauffement de Planck.
\end{solution}

\begin{solution}{exo:b2:anthropogenic-climate-change:4}
(a) Un pétagramme de carbone vaut $44.009/12.011 = 3.664$~Gt de dioxyde de carbone~: $0.45 \times 3.664 =
\qty{1.65}{\celsius}$ par \qty{1000}{PgC}, et l'intervalle devient de 0.989 à 2.31. (b) $2390 \times
12.011/44.009 = \qty{652}{PgC}$. (c) $0.45 \times 2.39 = \qty{1.08}{\celsius}$.
\end{solution}

\begin{solution}{exo:b2:anthropogenic-climate-change:5}
(a) $g = 159/398 = 0.400$~; $T = (398/\num{5.670e-8})^{1/4} = \qty{289}{K}$. (b) $-0.6 \times 4
\times \num{5.670e-8} \times 288^3 = \qty{-3.25}{W.m^{-2}.K^{-1}}$~; à \qty{298}{K}, $-3.60$.
(c) Moins sensible~: la réponse croît comme $T^3$, si bien qu'une planète plus chaude évacue
plus facilement chaque degré supplémentaire, les autres rétroactions étant égales par ailleurs.
\end{solution}

\begin{solution}{exo:b2:anthropogenic-climate-change:6}
(a) $1.85 - 0.50 = \qty{1.35}{W.m^{-2}.\celsius^{-1}}$. (b) $\sqrt{0.32^2 + 0.20^2} = 0.38$. (c) La somme est
bien mieux connue que chacun des deux termes~: leurs erreurs sont fortement anticorrélées. Là où la haute
troposphère se réchauffe davantage, elle contient aussi davantage de vapeur, si bien qu'une \omterm{def:b2:anthropogenic-climate-change:feedbacks}{rétroaction du
gradient thermique} négative plus forte va de pair avec une \omterm{def:b2:anthropogenic-climate-change:feedbacks}{rétroaction de la vapeur d'eau} positive plus forte.
\end{solution}

\begin{solution}{exo:b2:anthropogenic-climate-change:7}
(a) $\lambda = (0.59 - 2.20)/1.03 = \qty{-1.56}{W.m^{-2}.\celsius^{-1}}$~; $3.93/1.56 =
\qty{2.52}{\celsius}$. (b) $3.93/(1.56 - 0.5) = 3.93/1.06 = \qty{3.71}{\celsius}$. (c) $3.93/1.46 =
\qty{2.69}{\celsius}$. L'évaluation, avec des distributions complètes, trouve 3.5 et \qty{2.7}{\celsius}.
\end{solution}

\begin{solution}{exo:b2:anthropogenic-climate-change:8}
(a) $0.617/3682 = \qty{1.68e-4}{\celsius^{-1}}$. (b) $0.113 \times 0.396 = \qty{0.0447}{m}$,
soit \qty{44.7}{mm}, proche des \qty{47.5}{mm} du bilan. (c) $\num{2e-4} \times 0.2 \times 700 =
\qty{0.028}{m}$, soit \qty{28}{mm}.
\end{solution}

\begin{solution}{exo:b2:anthropogenic-climate-change:9}
(a) $\sum xy = 0.311$, $\sum x^2 = 0.300$~: $\beta = 1.04$. (b) L'intervalle à
\qty{95}{\percent} est $1.04 \pm 1.96 \times 0.08$, de 0.88 à 1.20~: il exclut 0, donc
la réponse est détectée, et il contient 1, donc le changement observé est compatible
avec le changement simulé. (c) L'intervalle de 0.34 à 0.66 exclut encore 0 (détecté),
mais exclut 1~: le modèle surestimerait la réponse d'environ un facteur deux.
\end{solution}

\begin{solution}{exo:b2:anthropogenic-climate-change:10}
(a) $\sum g^2 = 1.20$, $\sum a^2 = 0.18$, $\sum ga = -0.46$, $\sum gy = 0.74$, $\sum ay = -0.28$. Les
équations normales sont $1.20\,\beta_g - 0.46\,\beta_a = 0.74$ et $-0.46\,\beta_g + 0.18\,\beta_a
= -0.28$~; leur déterminant vaut $1.20 \times 0.18 - 0.46^2 = 0.0044$, et $\beta_g = (0.74 \times
0.18 - 0.46 \times 0.28)/0.0044 = 1.00$, $\beta_a = (-1.20 \times 0.28 + 0.46 \times 0.74)/0.0044 =
1.00$. (b) Gaz à effet de serre $+\qty{0.80}{\celsius}$, aérosols $-\qty{0.30}{\celsius}$. (c) Le
déterminant est petit parce que $g$ et $a$ sont presque proportionnels (leur corrélation vaut
$-0.99$)~: un peu de bruit sur $y$ peut être absorbé en augmentant ensemble les deux facteurs
d'échelle, si bien que chacun est mal contraint, bien que leur effet combiné soit bien connu.
\end{solution}

\begin{solution}{exo:b2:anthropogenic-climate-change:11}
(a) $(0.93 - 0.1)/0.45 \times 1000 = \qty{1844}{Gt}$~; moins 210 et $26 \times 0.93 = 24$~:
\qty{1610}{Gt}. (b) $0.93 - x = 0.45 \times (1350 + 210 + 24)/1000 = 0.713$, donc $x =
\qty{0.22}{\celsius}$~: à peu près la borne haute des 0.1 à \qty{0.2}{\celsius} évalués. (c)
Le budget ne dépend du réchauffement hors CO$_2$ $x$ que par $-x/\text{TCRE}$~: chaque
\qty{0.1}{\celsius} coûte $0.1/0.45 \times 1000 = \qty{222}{Gt}$, quelle que soit la limite.
\end{solution}

\begin{solution}{exo:b2:anthropogenic-climate-change:12}
(a) $3.93/1.81 = \qty{2.17}{\celsius}$, $3.93/0.51 = \qty{7.71}{\celsius}$, $3.93/1.16 =
\qty{3.39}{\celsius}$. (b) Le côté bas a pour longueur $3.39 - 2.17 = \qty{1.22}{\celsius}$, le côté haut
$7.71 - 3.39 = \qty{4.32}{\celsius}$~: la fonction $1/|\lambda|$ est convexe, si bien que des pas égaux de
$\lambda$ vers zéro élèvent l'ECS de plus en plus. (c) Un $\lambda$ faiblement négatif donne une ECS très
grande, et de petites variations de $\lambda$ près de zéro sont difficiles à exclure~; une ECS basse exige
un $\lambda$ fortement négatif, que la \omterm{def:b2:anthropogenic-climate-change:feedbacks}{réponse de Planck} borne et que le réchauffement observé contredit.
\end{solution}

\begin{solution}{pb:b2:anthropogenic-climate-change:1}
\textbf{1.} $2.16 + 0.54 + 0.21 + 0.41 = \qty{3.32}{W.m^{-2}}$.
\textbf{2.} $-0.22 - 0.84 = \qty{-1.06}{W.m^{-2}}$.
\textbf{3.} $3.32 + 0.47 + 0.05 - 1.06 - 0.20 + 0.08 + 0.06 = \qty{2.72}{W.m^{-2}}$, le total
évalué.
\textbf{4.} $3.93\log_2(410/278) = \qty{2.20}{W.m^{-2}}$, à moins de \qty{2}{\percent} de 2.16.
\textbf{5.} $3.93\log_2(427.35/278) = \qty{2.44}{W.m^{-2}}$.
\textbf{6.} Le terme aérosols--nuages, $-1.45$ à $-0.25$, large de \qty{1.20}{W.m^{-2}}, contre 1.52 pour
le total~: il porte l'essentiel de l'incertitude du total parce qu'il est grand
et mal contraint.
\textbf{7.} $0.79 \times \num{5.10e14} \times 12 \times \num{3.156e7} = \qty{1.53e23}{J}$,
\qty{153}{ZJ}.
\textbf{8.} $0.91 \times 153 = \qty{139}{ZJ}$.
\textbf{9.} $\lambda = (0.59 - 2.20)/1.03 = \qty{-1.56}{W.m^{-2}.\celsius^{-1}}$.
\textbf{10.} $3.93/1.56 = \qty{2.52}{\celsius}$.
\textbf{11.} $3.93 \times 1.03/2.20 = \qty{1.84}{\celsius}$.
\textbf{12.} $3.93/1.16 = \qty{3.39}{\celsius}$, plus élevée de \qty{0.87}{\celsius}.
\textbf{13.} Le réchauffement jusqu'ici a été le plus fort là où les rétroactions sont les plus négatives~; à mesure que le motif évolue,
le paramètre de rétroaction devient moins négatif, si bien que la sensibilité à long terme
dépasse la sensibilité effective.
\textbf{14.} $0.45 \times 44.009/12.011 = \qty{1.65}{\celsius}$ par \qty{1000}{PgC}.
\textbf{15.} 0.989 à \qty{2.31}{\celsius} par \qty{1000}{PgC}.
\textbf{16.} $2390 \times 12.011/44.009 = \qty{652}{PgC}$.
\textbf{17.} $0.45 \times 2.39 = \qty{1.08}{\celsius}$.
\textbf{18.} Aux arrondis près, le dioxyde de carbone seul explique tout le réchauffement humain~:
le réchauffement dû aux autres gaz à effet de serre et le refroidissement dû aux aérosols se compensent presque.
\textbf{19.} $1.5 - 1.07 = \qty{0.43}{\celsius}$.
\textbf{20.} $(0.43 - 0.1)/0.45 \times 1000 = \qty{733}{Gt}$ de dioxyde de carbone.
\textbf{21.} $733 - 210 = \qty{523}{Gt}$.
\textbf{22.} $26 \times 0.43 = \qty{11.2}{Gt}$~; $523 - 11.2 = \qty{512}{Gt}$.
\textbf{23.} Il baisse de $0.1/0.45 \times 1000 = \qty{222}{Gt}$, pour atteindre \qty{290}{Gt}~: les
\qty{220}{Gt} dont l'évaluation dit que les choix sur les autres facteurs déplacent le budget.
\textbf{24.} Calculé à partir de la TCRE, le \omterm{def:b2:anthropogenic-climate-change:budget}{budget carbone restant} pour \qty{1.5}{\celsius} à partir du
début de 2020 est de \textbf{\qty[text-series-to-math=true]{512}{Gt}} de dioxyde de carbone, à \qty{3}{\percent} près
du \qty{500}{Gt} évalué avec une probabilité de \qty{50}{\percent}.
\end{solution}
